\scanpage{202}\mspage{14\deleted{\unsure{6}}7}
\deleted{\unsure{tout|pour} $Y$ \unsure{de|dans} $\mathrm{Cat}_{/S}$, $X\times_S Y\longrightarrow X'\times_S Y$ est dans $W_S^d$.}

On a ainsi prouvé que les foncteurs
\diagram{\begin{tikzcd}[ampersand replacement=\&,column sep=large]
 (*) \& \mathrm{HOT}_{\unsure{W|w}}(S') \arrow[r,shift left=2pt,"f_!"] \& \mathrm{HOT}_W(S) \arrow[l,shift left=2pt,"f^*"]
\end{tikzcd}}
sont adjoints l'un de l'autre, avec des morphismes d'adjonction
\[
 \left\{\begin{aligned}
 \mathrm{id}&\xrightarrow{\alpha} f^*f_!\\
 \mathrm{id}&\xleftarrow{\beta} f_!f^*
 \end{aligned}\right.
\]
précisés plus haut. Je présume que la même démonstration, en travaillant avec les \deleted{\unsure{relatifs}} localiseurs \inserted{grossiers} $W_{\unsure{S\to|S'|S}},W_{S'}$ au lieu de $W_S^d,W_{S'}^d$, donnerait que les foncteurs correspondants
\diagram{\begin{tikzcd}[ampersand replacement=\&,column sep=large]
 (**) \& \mathrm{HOT}_W^{\mathrm{lc}}(S') \arrow[r,shift left=2pt,"f_!^{\mathrm{lc}}"] \& \mathrm{HOT}_W^{\mathrm{lc}}(S) \arrow[l,shift left=2pt,"f^*_{\mathrm{lc}}"]
\end{tikzcd}}
sont adjoints l'un de l'autre, avec essentiellement ``les mêmes'' morphismes d'adjonction $\alpha,\beta$. Je n'ai pas envie de reprendre les raisonnements précédents, et préfère le \deleted{\unsure{prouver}} \uline{déduire} de l'adjonction dans la situation $(*)$. Cela me donnera l'occasion de mieux comprendre les traits formels essentiels de la situation. Je poserai pour simplifier
% archive sheet number: 201
